\documentclass[12pt, a4paper]{article} 
\usepackage[siunitx]{circuitikz}
\usepackage[utf8]{inputenc} 
\usepackage[spanish]{babel} 
    \usepackage{amsmath} 
    \usepackage{geometry} 
\geometry{margin=2.5cm}

\title{\textbf{Preinforme}}
\author{Maria José Álvarez Hernández, Cristian Alejandro Guarin Marin}
\date{\today}

\begin{document}

\maketitle

\section{Consulta}
Como grupo comprendimos que las transformaciones delta-estrella y 
estrella-delta son técnicas de simplificación de circuitos. Las utilizamos 
cuando nos encontramos con redes de resistencias que no están conectadas en 
configuraciones de serie ni paralelo.

En esencia estas transformaciones nos permiten reemplazar un arreglo de
tres resistencias en forma de Y por un arreglo equivalente con un nodo
central común, o viceversa. El objetivo de hacer este reemplazo es destrabar el circuito, logrando que las resistencias restantes queden claramente en serie 
o en paralelo, cosa que nos permitirá encontrar la $R_{eq}$. Cabe destacar
que al hacer estas transformaciones el voltaje y la corriente en los terminales
externos del circuito se mantienen exactamente iguales.

\subsection{Transformaciones delta-estrella}
La regla general es que la resistencia de la estrella en un nodo dado es
igual al producto de las dos resistencias adyacentes de la delta, dividido 
por la suma de las tres resistencias de la delta.

\vspace{0.5cm}
$$R_A = \frac{R_{AB}R_{CA}}{R_{AB}+R_{BC}+R_{CA}}$$
\vspace{0.5cm}
$$R_B = \frac{R_{AB}R_{BC}}{R_{AB}+R_{BC}+R_{CA}}$$
\vspace{0.5cm}
$$R_C = \frac{R_{CA}R_{BC}}{R_{AB}+R_{BC}+R_{CA}}$$

\vspace{0.5cm}

\subsection{Transformaciones estrella-delta}
La regla general es que la resistencia de la delta conectada entre dos nodos
es igual a la suma de todos los productos cruzados posibles de las resistencias
de la estrella, dividida por la resistencia de la estrella conectada al nodo 
opuesto.

\vspace{0.5cm}
$$R_{AB} = \frac{R_AR_B+R_BR_C+R_CR_A}{R_C}$$
\vspace{0.5cm}
$$R_{BC} = \frac{R_AR_B+R_BR_C+R_CR_A}{R_A}$$
\vspace{0.5cm}
$$R_{CA} = \frac{R_AR_B+R_BR_C+R_CR_A}{R_B}$$

\section{Cálculos}

%Circuito 1------------------------------------------------------
\subsection{Circuito 1}

\vspace{0.8cm}

\begin{center}
\begin{circuitikz}
    %Fuente de voltaje en el extremo izquierdo
    \draw (0,0) to[battery1, invert, l=$\text{8V}$] (0,4);
    
    %Rieles superior e inferior
    \draw (0,4) to[short, i=$I$] (9,4);
    \draw (0,0) to[short] (9,0);

    %Primera resistencia bajando en x=3
    \draw (3,4) to[R, i=$I_A$, l=$1.5k\varOmega$] (3,0);

    %Segunda resistencia bajando en x=6
    \draw (6,4) to[R, i=$I_B$, l=$2.2k\varOmega$] (6,0);

    %Tercera resistencia bajando en x=9
    \draw (9,4) to[R, i=$I_C$, l=$3.3k\varOmega$] (9,0);
\end{circuitikz}
\end{center}

\vspace{0.8cm}

%Tabla
\begin{table}[h]
\centering
\begin{tabular}{|c|c|c|c|}
\hline
\textbf{$R_{eq1}$} & \textbf{$I_A$} & \textbf{$I_B$} & \textbf{$I_C$} \\
\hline
$702.12\varOmega$ & $5.3mA$ & $3.6mA$ & $2.4mA$ \\
\hline
\end{tabular}
\label{tab:tabla_2x4}
\end{table}

%Circuito 2 -------------------------------------------------------------
\subsection{Circuito 2}

\vspace{0.8cm}

\begin{center}
\begin{circuitikz}
    %Fuente de voltaje en el extremo izquierdo
    \draw (0,0) to[battery1, invert, l=$8V$] (0,4);

    %Resistencia superior
    \draw (0,4) to[R, i=2mA, l=$R_a1k\varOmega$] (4,4);

    %Resistencia derecha
    \draw (4,4) to[R, l=$R_b2k\varOmega$] (4,0);

    %Resistencia inferior
    \draw (4,0) to[R, l=$R_c1k\varOmega$] (0,0);
\end{circuitikz}
\end{center}

%Tabla
\begin{table}[h]
    \centering
    \begin{tabular}{|c|c|c|c|c|}
        \hline 
        $\mathbf{R_{eq2}}$ & $\mathbf{i}$ & $\mathbf{V_{R_a}}$ & $\mathbf{V_{R_b}}$ & $\mathbf{V_{R_c}}$ \\
        \hline
        $4k\varOmega$ & $2mA$ & $2V$ & $4V$ & $2V$ \\
        \hline
    \end{tabular}
\end{table}

%Circuito 3----------------------------------------------------------------------------------------------------
\subsection{Circuito 3}

\vspace{0.8cm}

\begin{center}
\begin{circuitikz}
    %Fuente alienada a la izquierda
    \draw (0,0) to[battery1, invert, l=$8V$] (0,10);

    %Rieles superior e inferior
    \draw (0,10) to[short, i=$i$] (10,10);
    \draw (0,0) to[short] (10,0);

    %Primer paralelo
    \draw (2,10) to[R, l=$R_a$] (2,5);
    \draw (2,5) to[R, l=$R_d$] (2,0);

    %Horizontal
    \draw (2,5) to[R, l=$R_c = 1k\Omega$] (10,5);

    %Segundo paralelo
    \draw (10,10) to[R, l=$R_b = 1k\Omega$] (10,5);
    \draw (10,5) to[R, l=$R_e = 1k\Omega$] (10,0);
\end{circuitikz}
\end{center}

%Tabla
\begin{table}[h]
    \centering
    \begin{tabular}{|c|c|c|c|c|c|c|}
        \hline
        $\mathbf{R_{eq3}}$ & $\mathbf{R_a}$ & $\mathbf{R_b}$ & $\mathbf{R_c}$ & $\mathbf{R_d}$ & $\mathbf{R_e}$ & $\mathbf{I}$\\
        \hline
        $1k\varOmega$ & $1k\varOmega$ & $1k\varOmega$ & $1k\varOmega$ & $1k\varOmega$ & $1k\varOmega$ & $9mA$ \\
        \hline

        
    \end{tabular}
\end{table}

\end{document}